\subsection{首批应用}

设 A 为诺特环，V = Spec(A) 为其谱。在本节中，我们把 V 中任何形如 \footnote{回忆，\(V(x)\) 由 \(A\) 的包含 \(x\) 的素理想组成。} V(x) 且 \(x \in A\) 的子集称为超曲面。

命题 4. --- 设 X 是 V 的闭子集，\(H_{1}, \ldots, H_{m}\) 是 V 中的超曲面。置 \(X' = X \cap H_{1} \cap \ldots \cap H_{m}\)。

\begin{enumerate}
\def\labelenumi{\alph{enumi})}
\tightlist
\item
  对于包含于 X′ 的 V 的每个闭子集 Y，有
\end{enumerate}

\[
\operatorname{codim} (\mathrm{Y}, \mathrm{X} ^ {\prime}) \geqslant \operatorname{codim} (\mathrm{Y}, \mathrm{X}) - m.
\]

\begin{enumerate}
\def\labelenumi{\alph{enumi})}
\setcounter{enumi}{1}
\item
  对于 X' 的每个不可约分支 Z，有 codim(Z, X) ≤ m。若 X' 非空，则 codim(X', X) ≤ m。
\item
  若 Z 是包含于 X 的 V 的闭不可约子集，且 \(\text{codim}(Z, X) \leqslant m\)，则存在超曲面 \(H_{1}^{\prime}, \ldots, H_{m}^{\prime}\)，使 Z 成为 \(X \cap H_{1}^{\prime} \cap \ldots \cap H_{m}^{\prime}\) 的不可约分支。
\end{enumerate}

设 \(\mathfrak{a}\) 是 A 的理想，\(x_{1},\ldots,x_{m}\) 是 A 中的元素，使得 \(\mathrm{X}=\mathrm{V}(\mathfrak{a})\)，且 \(\mathrm{H}_{i}=\mathrm{V}(x_{i})\)（\(1\leqslant i\leqslant m\)）。设 Z 是包含于 X 的 V 的闭不可约子集；则存在包含 \(\mathfrak{a}\) 的 A 的素理想 p，使 \(\mathrm{Z}=\mathrm{V}(\mathfrak{p})\)。

首先假设 Z 包含于 X'，并记 \(\xi_{i}\) 为 \(x_{i}\) 在诺特局部环 \(\mathbf{B} = \mathbf{A}_{\mathfrak{p}} / \mathfrak{a}\mathbf{A}_{\mathfrak{p}}\) 中的像。由 § 1, no. 3 的 prop. 7, b)，有

\[
\operatorname{codim} (Z, X) = \dim (B), \quad \operatorname{codim} (Z, X ^ {\prime}) = \dim (B / (\xi_ {1} B + \dots + \xi_ {m} B)).
\]

因此，由 no. 1 的 prop. 2，有

\[
\operatorname{codim} (Z, X ^ {\prime}) \geqslant \operatorname{codim} (Z, X) - m.\tag{11}
\]

若 Z 是 X' 的不可约分支，则 codim(Z, X') = 0，因而 codim(Z, X) ≤ m；这证明了 b)。对 Y 的不可约分支 Z 的集合，在 (11) 的两边分别取下界，即可证明 a)。

反之，假设 \(\operatorname{codim}(\mathbf{Z},\mathbf{X})\leqslant m\)，即 \(\dim (\mathbf{B})\leqslant m\)。由于 \(\mathbf{A}_{\mathfrak{p}}\) 中的每个元素都是 \(\mathbf{A}_{\mathfrak{p}}\) 中的可逆元素与 A 中某个元素的像之积，No. 2 的附论证明存在 A 中的元素 \(x_1^{\prime},\ldots ,x_m^\prime\)，它们在 B 中的像生成 B 的一个定义理想。置 \(\mathrm{H}_i^{\prime} = \mathrm{V}(x_i^{\prime})\)，其中 \(1\leqslant i\leqslant m\)。显然，Z 是 \(\mathbf{X}\cap \mathbf{H}_1^{\prime}\cap \dots \cap \mathbf{H}_m^{\prime}\) 的不可约分支。

推论 1. --- 设 H 是 V 中的非空超曲面。H 在 V 中的余维等于 0 或 1。

codim(H, V) = 1 当且仅当 H 不包含 V 的任何不可约分支。在这种情况下，H 的所有不可约分支在 V 中的余维都是 1。

对于 H 的每个不可约分支 Z，有

\[
0 \leqslant \operatorname{codim} (Z, V) \leqslant 1
\]

根据 prop. 4, b)，并且 codim(Z, V) = 0 当且仅当 Z 是 V 的不可约分支。根据定义，

\[
\operatorname{codim} (\mathrm{H}, \mathrm{V}) = \inf _ {\mathrm{Z}} \operatorname{codim} (\mathrm{Z}, \mathrm{V})
\]

其中 Z 遍历 H 的不可约分支的集合。推论 1 立即由这些备注得到。

推论 2. --- 设 X 是 V 的闭不可约子集，H 是 V 中的超曲面。只有以下三种情形：

\begin{enumerate}
\def\labelenumi{\arabic{enumi})}
\item
  有 \(\mathbf{X}\subset \mathbf{H}\)
\item
  集合 \(\mathbf{X} \cap \mathbf{H}\) 非空，且它的每个不可约分支 \(Z\) 都满足 \(\operatorname{codim}(\mathbf{Z}, \mathbf{X}) = 1\)；
\item
  集合 X ∩ H 为空。
\end{enumerate}

假设 \(X' = X \cap H\) 非空且不同于 \(X\)；每个不可约分支 \(Z\) 属于 \(X'\)，都不同于 \(X\)，且满足 codim \((Z, X) \leqslant 1\)，由 prop. 4, b)。推论 2 立即由此得到。

推论 3. --- 若 A 是唯一分解环（VII, § 3, no. 3, prop. 2），则 A 的高度为 1 的素理想正是由 A 的极元生成的主理想。若 A 还是局部环，则对于 A 的每个高度为 1 的素理想 p，都有 dim(A/p) = dim(A) - 1。

设 \(x\) 是 A 的极元。则 \(Ax\) 是素理想，因为 \(x\) 是极元；它的高度为 1，因为 A 是整环（no. 1, prop. 1）。设 \(p\) 是 A 的高度为 1 的素理想。则 \(V(p)\) 是超曲面 \(V(Ax)\) 的一个不可约分支，对于适当的 \(x\)（prop. 4, c)）。设 \(x = \prod_{i} y_{i}^{n_i}\) 是把 \(x\) 分解为极元之积的分解，

其中 \(y_{i}\) 与 \(y_{j}\) 当 \(i \neq j\) 时互素。\(V(Ax)\) 的不可约分支是 \(V(Ay_{i})\)。因此对于适当的 i，有 \(p = Ay_{i}\)。最后一个断言由 no. 1 的 prop. 2 的 cor. 2 得到。

备注。--- 1) 假设 A 是维数为 d 的局部诺特整环；令 x 为 \(m_{A}\) 中的非零元素，并令 \(H = V(x)\)。由 prop. 4 的 cor. 2，H 的每个不可约分支在 X 中的余维为 1，因而维数 \(\leqslant d - 1\)（§ 1, no. 2, prop. 3）。由 no. 1 的 prop. 2 的 cor. 2，H 的维数为 d - 1，因而这些分支中有一个维数为 d - 1；若 A 是串链环，则它们全都如此。然而，一般而言，H 可能有一个维数 \textless{} d - 1 的不可约分支（cf. p. 87, exercise 7）。

\begin{enumerate}
\def\labelenumi{\arabic{enumi})}
\setcounter{enumi}{1}
\item
  设 \(x_1, \ldots, x_n\) 是 A 中的元素，并置 \(\mathrm{H}_i = \mathrm{V}(x_i)\)，其中 \(1 \leqslant i \leqslant n\)。假设存在一个支集为 \(\mathrm{V} = \operatorname{Spec}(\mathrm{A})\) 的有限生成 A-模 M，使 \((x_1, \ldots, x_n)\) 是 M 的完全截切序列。则 \(\mathrm{H}_1 \cap \ldots \cap \mathrm{H}_n\) 的每个不可约分支在 V 中的余维都是 \(n\)：这由 No. 2 的 Prop. 3 的推论容易得到。
\item
  若诺特环 A 的理想 a 由 m 个元素生成，则 ht(a) ≤ m。这立即由 prop. 4 得到。
\end{enumerate}

命题 5. --- 设 A 为诺特环，\(\mathfrak{p} \subset \mathfrak{q}\) 是 A 的一个非饱和素理想链。由 A 的素理想 \(\mathfrak{r}\) 中满足 \(\mathfrak{p} \subset \mathfrak{r} \subset \mathfrak{q}\) 者组成的集合 E 是无限链。我们有 \(\bigcup \mathfrak{r} = \mathfrak{q}\) 且 \(\bigcap \mathfrak{r} = \mathfrak{p}\)。

将 A 替换为 A/p 后，可归结为 p = \{0\} 的情形。

由 Section 1 的引理 1，有 q = \(\bigcup_{r \in E} r\)，而 II, § 1, no. 1 的 prop. 2 表明 E 是无限的。

令 \(y \neq 0\) 为 \(\bigcap_{r \in E} r\) 中的元素。q 的高度是有限的（\(n^{\circ}1\)，prop. 2 的 cor. 1），并且根据假设有 \(ht(q) \geqslant 2\)。因此存在一个高度为 2 的素理想 \(q' \subset q\)。将证明的第一部分应用于 \(q'\)，可知集合 \(E'\)（由包含于 \(q'\) 的高度为 1 的素理想组成）是无限的；根据假设，这些理想中的每一个都包含 y。现在，根据 n° 1 的 prop. 2 的 cor. 2，诺特局部环 \(B = A_{q'} / yA_{q'}\) 的维数为 1。因此，对每个 \(r \in E'\)，B 的素理想 \(r/yA_{q'}\) 是极小的；于是诺特环 B 有无穷多个极小素理想，这是矛盾的（II, § 4, n° 3, cor. 3 of prop. 14）。所以有 \(\bigcap_{r \in E} r = \{0\}\)。

命题 6.--- 设 A 是维数 \(\geqslant 2\) 的诺特环，\(h\) 是满足 \(0 < h < \dim(A)\) 的整数。

\begin{enumerate}
\def\labelenumi{\alph{enumi})}
\item
  A 有无穷多个高度为 h 的素理想。
\item
  若 A 的维数有限，则 A 有无穷多个素理想 p，使 \(\operatorname{ht}(\mathfrak{p}) = h\) 且 \(\dim(\mathbf{A}/\mathfrak{p}) = \dim(\mathbf{A}) - h\)。
\end{enumerate}

由于 A 的维数是 A 的素理想的高度的上界（§ 1, n° 3, prop. 8），由于 A 的每个素理想都有有限高度（n° 1, prop. 2 的 cor. 1），又由于 \(h < \dim(A)\)，存在整数 \(n > h\) 以及素理想 \(p\)（高度为 \(n\)），因而存在素理想链 \(p_0 \subset \ldots \subset p_n = p\)，其长度为 \(n\)。有 \(ht(p_i) = i\)，其中 \(0 \leqslant i \leqslant n\)，所以有 \(ht(r) = h\)，对于 A 中每个素理想 \(r\)，满足 \(p_{h-1} \subset r \subset p_{h+1}\)。根据 prop. 5，这些理想组成的集合 E 是无限的，故得 \(a\) ).

若 A 的维数有限，则在上述论证中可以设 \(n = \dim(A)\)。对于每个素理想 \(r \in E\)，有 \(\operatorname{ht}(r) = h\)，从而 \(\dim(A/r) \leqslant n - h\)；又由于 \(r \subset p_{h+1} \subset \ldots \subset p_n\) 是长度为 \(n - h\) 的链，有 \(\dim(A/r) = n - h\)。

存在维数为 2 且只有一个高度为 1 的素理想的非诺特整环，例如高度为 2 的赋值环（VI, § 4, n° 4, prop. 5）。

